% ---------------------------------------------------------------------
\subsection{Lucier's room (\fn{Lucier.cpp})}
\label{sec:math-lucier}
The Lucier mode (\S\ref{sec:lucier}) plays the organ through the response of $n$ passes of
\enquote{loudspeaker -- room -- microphone}. Everything is computed in the frequency domain once per change, on
a background thread, and the result is played by the same partitioned convolution as the room tail.

\paragraph{The room's response (\fn{lucierFromRoom}).} For the simulated room the single-pass response is built from
two parts. Above a crossover frequency it is the diffuse field of \S\ref{sec:math-tail} (two-band decaying noise,
here with the Eyring times and its own seed). Below it, where a room has few, separated resonances, it is the sum of
the room's eigenmodes. For a rectangular room with rigid walls the modes are
\begin{equation}
  f_{n_x n_y n_z} = \frac c2\sqrt{\Bigl(\frac{n_x}{L_x}\Bigr)^2 + \Bigl(\frac{n_y}{L_y}\Bigr)^2 + \Bigl(\frac{n_z}{L_z}\Bigr)^2},
\end{equation}
with the mode shapes $\psi = \prod_i\cos(n_i\pi x_i/L_i)$. The transfer function from a source at $\mathbf s$ (the
organ, \SI{6}{m} in front, $10^\circ$ up) to the listener $\mathbf r$ is the modal sum
\begin{equation}
  M(\omega) = \sum_{\mathbf n} \varepsilon_{\mathbf n}\,\psi_{\mathbf n}(\mathbf s)\,\psi_{\mathbf n}(\mathbf r)\,
  \frac{i\omega}{\omega_{\mathbf n}^2 - \omega^2 + 2i\delta\omega},
\end{equation}
a sum of resonators, each driven according to how strongly source and listener couple to its shape, with the
normalisation $\varepsilon = \prod_i(2-\delta_{n_i0})$ and the damping $\delta = 6.9078/T_\mathrm{lo}$ (the amplitude
decay rate that gives $T_\mathrm{lo}$; the half-power bandwidth of each mode is then $\delta/\pi$). Each mode is only
evaluated within $\pm\max(\SI{3}{Hz}, 12\delta/2\pi)$ of its frequency, where its contribution is not negligible. The
crossover is Schroeder's frequency $f_S = 2000\sqrt{T/V}$, doubled and clamped to 60--\SI{400}{Hz}, and lowered until
the expected number of modes below $1.3f_c$,
\begin{equation}
  N(f) \approx \frac{4\pi V}{3}\Bigl(\frac fc\Bigr)^3 + \frac{\pi S}{8}\Bigl(\frac fc\Bigr)^2 + \frac{L_x+L_y+L_z}{2}\,\frac fc
\end{equation}
(the volume, surface and edge terms of the mode count; the textbook surface term is $\pi S/4$, so the code's estimate
is slightly low and a few more modes may be computed than the limit says), is at most 4000. The two parts are joined with a raised-cosine
crossfade from $0.8f_c$ to $1.25f_c$, the modes scaled so that their mean energy in $[f_c/2, f_c]$ equals that of the
diffuse field in $[f_c, 2f_c]$. A polygon room uses the modes of a rectangle with the same floor area (the exact modes
of an arbitrary plan would need a wave solver). A loaded impulse response is used as it is, from \SI{2}{ms} before its
direct sound (first channel only).

\paragraph{One pass and $n$ passes (\fn{buildLucier}).} Every pass also goes through a loudspeaker and a microphone,
modelled as a band-pass from about 60\,Hz to \SI{10}{kHz},
\begin{equation}
  |T(f)| = \frac{q^2}{\sqrt{1+q^4}}\cdot\frac{1}{\sqrt{1+(f/10^4)^2}},\qquad q = (f/60)^2,
\end{equation}
(a fourth-order high-pass and a first-order low-pass). The response is damped by $e^{-t/\tau}$, with
$\tau = \min(\tau_\mathrm{user},\, 2.5/(n-1))$ so that later generations stay within a few seconds, multiplied by $T$
and normalised to unit RMS between 60\,Hz and \SI{10}{kHz}; call the result $H(f)$. Playing a sound $n$ times through
the room multiplies its spectrum by $H^n$. For an integer $n$ this is computed exactly as
$|H|^n e^{in\arg H}$ (no phase unwrapping is needed, since $n\arg H$ is right modulo $2\pi$). Each power is scaled by
\begin{equation}
  g_n = \bigl(\max|H^n|\cdot\operatorname{rms}|H^n|\bigr)^{-1/2},
\end{equation}
the geometric mean between \enquote{same peak} and \enquote{same energy}: normalising to the peak would let the rest
vanish, normalising to the energy would let the one resonance that survives explode. A fractional generation
$n_0+\phi$ crossfades the spectra of $n_0$ and $n_0+1$ linearly. As $n$ grows, $|H|^n$ is dominated more and more by
the largest peaks of $|H|$, which is the process of Lucier's piece: the room's resonant frequencies remain, everything
else fades. The display shows the strongest local maxima of $|G|$ that lie at least a sixth of an octave apart.

\paragraph{The piece (\fn{lucierPass}).} For \ui{Stück}, a whole recording $x$ is convolved with the single-pass response
$h$ again and again, each time by FFT overlap-add with blocks of $B = N-\ell_h+1$ samples ($N = 2^k\ge 2\ell_h$), so
that every block's linear convolution of length $B+\ell_h-1 = N$ fits into one transform without wrapping around.

% ---------------------------------------------------------------------
\subsection{The output stage (\fn{MasterBus.h})}
\label{sec:math-master}
\paragraph{Equaliser.} Three biquads after Bristow-Johnson's cookbook \citep{rbj}, the same on every channel (an Ambisonic
field stays a field: a filter that is the same on all channels commutes with every rotation and decoder). With
$A = 10^{G/40}$, $\omega = 2\pi f/f_s$:
\begin{itemize}[nosep]
  \item bell: $\alpha = \sin\omega/(2Q)$, $H(z) = \dfrac{(1+\alpha A) - 2\cos\omega\,z^{-1} + (1-\alpha A)z^{-2}}{(1+\alpha/A) - 2\cos\omega\,z^{-1} + (1-\alpha/A)z^{-2}}$,
    which has the gain $A^2 = 10^{G/20}$ at $f$ and 1 far from it;
  \item shelves with slope $S=1$: $\alpha = \tfrac{\sin\omega}{2}\sqrt2$, $\beta = 2\sqrt A\,\alpha$, $c = \cos\omega$, and
    the low shelf
    \begin{align*}
      b_0 &= A\bigl[(A+1)-(A-1)c+\beta\bigr], & a_0 &= (A+1)+(A-1)c+\beta,\\
      b_1 &= 2A\bigl[(A-1)-(A+1)c\bigr], & a_1 &= -2\bigl[(A-1)+(A+1)c\bigr],\\
      b_2 &= A\bigl[(A+1)-(A-1)c-\beta\bigr], & a_2 &= (A+1)+(A-1)c-\beta,
    \end{align*}
    with the gain $A^2$ at DC and 1 at Nyquist; the high shelf has the signs of the $(A-1)c$ terms and of $b_1, a_1$
    reversed. All coefficients are divided by $a_0$.
\end{itemize}
Each biquad runs in transposed direct form II,
$y = b_0x + z_1$, $z_1 \leftarrow b_1x - a_1y + z_2$, $z_2 \leftarrow b_2x - a_2y$, which has the best numerical behaviour
of the direct forms in floating point; states below $10^{-20}$ are flushed to zero so silence does not produce
denormal numbers. A band at $0$\,dB is skipped. The curve on screen is $|H(e^{i\omega})|$ of the product, evaluated
directly from the coefficients.

\paragraph{Compressor.} The detector takes the largest $|x_c|$ of all channels, in dB, and follows it with separate
attack and release coefficients, $E \leftarrow L + a\,(E - L)$, $a = e^{-1/(\tau f_s)}$ with $\tau$ the attack time when
$L > E$ and the release time otherwise. The gain reduction is the soft-knee curve of width $W = \SI{6}{dB}$ with the
threshold $T$ and ratio $R$ (over $= E - T$, slope $s = 1/R - 1$):
\begin{equation}
  G_r = \begin{cases} 0, & 2\,\text{over} < -W\\[2pt] s\,\dfrac{(\text{over} + W/2)^2}{2W}, & |2\,\text{over}|\le W\\[6pt]
  s\cdot\text{over}, & 2\,\text{over} > W\end{cases}
\end{equation}
a quadratic that joins the two straight lines with matching value and slope at both ends of the knee. The gain itself
glides with \SI{2}{ms} and is applied, with the make-up gain, to every channel alike. Because all channels share
one gain, the image does not move when the compressor works: a loud pipe on the left does not push the rest of the
field to the right.

% ---------------------------------------------------------------------
\subsection{Modulation (\fn{Modulation.h})}
\label{sec:math-mod}
\paragraph{Hashing (\fn{hash01}).} Many random values must be the same every time they are asked for (a step of an
LFO, a pipe's offset, a pass of a lane). They come from an integer hash, not a generator with state:
\begin{gather}
  a \leftarrow a\oplus(a\gg16),\quad a \leftarrow a\cdot\mathtt{0x7feb352d},\quad a \leftarrow a\oplus(a\gg15),\\
  a \leftarrow a\cdot\mathtt{0x846ca68b},\quad a \leftarrow a\oplus(a\gg16),
\end{gather}
returning the low 24 bits divided by $2^{24}-1$. Each multiply spreads every input bit over the higher bits, each
shift-and-xor folds the high bits back, so neighbouring inputs give unrelated outputs (a well-mixing finaliser of the
\enquote{lowbias32} family).

\paragraph{LFO shapes (\fn{lfoValue}).} With the phase $\Phi$ in cycles and $f = \Phi - \lfloor\Phi\rfloor$, all shapes
start at 0 and rise first, aligned with the sine:
\begin{gather}
  \text{sine } \sin 2\pi\Phi,\qquad
  \text{triangle } \begin{cases}4f, & f<\tfrac14\\2-4f, & f<\tfrac34\\4f-4, & \text{else}\end{cases}\qquad
  \text{saw } \begin{cases}2f, & f<\tfrac12\\2f-2, & \text{else}\end{cases}\qquad
  \text{square } \pm1,\\
  \text{steps } r_{\lfloor\Phi\rfloor},\qquad
  \text{smooth } r_{\lfloor\Phi\rfloor} + \bigl(r_{\lfloor\Phi\rfloor+1}-r_{\lfloor\Phi\rfloor}\bigr)(3f^2-2f^3),
\end{gather}
where $r_i = 2\,\mathrm{hash01}(i\cdot2654435761 \oplus \text{seed})-1$ is one random value per cycle (2654435761 is
Knuth's multiplicative constant $\approx 2^{32}/\varphi$). The smooth shape interpolates between them with the smoothstep,
which has zero slope at the cycle boundaries, so it has no corners. Because the value depends only on the phase, a
host that jumps or loops gets the same LFO back. The phase is integrated from the rate (or from the host's tempo and the
division in quarter notes per cycle), and each pipe adds an offset $\sigma\kappa_k$ to it (\ui{Versatz}): along the
keyboard $\kappa_k = k/41$, or random per pipe $\kappa_k = \mathrm{hash01}(131s + 7k + 1)$, so at $\sigma = 1$ one cycle
is spread over the whole stop and the modulation runs through the pipes like a wave.

\paragraph{Envelopes (\fn{EnvState}).} Delay, attack, decay, sustain, release, advanced every control block by $\Delta t$.
The attack is linear, $\ell \mathrel{+}= \Delta t/T_A$, since a linear rise sounds even for a level that drives other
parameters; decay and release are exponential towards their targets,
\begin{equation}
  \ell \leftarrow S + (\ell - S)\,e^{-4\Delta t/T_D},\qquad \ell \leftarrow \ell\, e^{-4\Delta t/T_R},
\end{equation}
so that after the nominal time $T$ the remaining distance is $e^{-4}\approx 2\,\%$ (\SI{-35}{dB}); the stage ends when the
distance falls below $10^{-3}$ (decay) or the level below $5\cdot10^{-4}$ (release). A new trigger restarts from the
current level, so retriggering does not click.

\paragraph{Routing.} Each of the 24 routings adds $a\cdot s'$ to its target, where $s'$ is the source, mapped to $[0,1]$
by $(s+1)/2$ when the routing is unipolar and the source bipolar. For a per-pipe target the sum of all routings $m$ acts
on each pipe with its own source values (its LFO offset, its envelope, its key, its velocity, its MPE pressure), with the
scales of \S\ref{sec:math-voice} ($180^\circ m$, $2^{3m}$, $36m$ dB, \dots). For a parameter of the plugin it is added
to the parameter's normalised value.

% ---------------------------------------------------------------------
\subsection{Automation over time (\fn{Automation.h})}
\label{sec:math-verlauf}
The value of a lane is a pure function of the clock time $t$ in beats, so seeking, host loops and bounces give exactly
the same result as playing through.

\paragraph{Lane time (\fn{autoLocalTime}).} With the lane's offset $t_0$, speed $v$ and length $\Lambda$,
$\tau = (t-t_0)\,v$; a looping lane takes $p = \lfloor\tau/\Lambda\rfloor$ as the number of the pass and
$\tau \leftarrow \tau - p\Lambda$; a reversed lane reads $\Lambda - \tau$.

\paragraph{Segments (\fn{autoSegment}).} Between point $a$ at $t_a$ and point $b$ at $t_b$, with $u = (\tau-t_a)/(t_b-t_a)$
and the bend $\beta\in[-1,1]$ of $a$,
\begin{equation}
  B(u) = u^{\,2^{-3\beta}},
\end{equation}
which runs from $u^8$ (slow start, fast end) through the straight line at $\beta=0$ to $u^{1/8}$ (fast start). The
shapes are
\begin{itemize}[nosep]
  \item \ui{Kurve}: $v_a + (v_b-v_a)B(u)$; \ui{Stufe}: $v_a$ until the next point;
  \item \ui{Treppe}: $n = \max(1,\mathrm{round}(\rho\,(t_b-t_a)))$ steps, $u\leftarrow\min(1, \lfloor nu\rfloor/(n-1))$, then the curve;
  \item \ui{Welle}: the curve plus $\tfrac D2\sin(2\pi\rho(\tau-t_a))$ with the rate $\rho$ in cycles per beat and the depth $D$;
  \item \ui{Zittern}: the curve plus $\tfrac D2\,\nu(\rho(\tau-t_a))$, with value noise
    $\nu(x) = r_{\lfloor x\rfloor} + (r_{\lfloor x\rfloor+1}-r_{\lfloor x\rfloor})\bigl(\tfrac12-\tfrac12\cos\pi\{x\}\bigr)$ seeded by the segment.
\end{itemize}
The result is clamped to $[0,1]$ and, for a switch or a choice with $k$ steps, rounded to $\mathrm{round}(v(k-1))/(k-1)$.

\paragraph{Variation from pass to pass.} In pass $p$ a point with \ui{Streuung} $\sigma$ takes the value
$v + \sigma(2\,\mathrm{dice}(p,i,\cdot)-1)$; the segment from point $i$ moves only if $p \bmod n = n-1$ for \ui{jedes $n$-te
Mal} and $\mathrm{dice}(p,i,\cdot) < $ \ui{Chance}, otherwise it holds its start value. The dice
$\mathrm{dice}(p,i,s) = \mathrm{hash01}\bigl(2654435761p \oplus 40503i \oplus 97\,\mathrm{seed}\oplus s\bigr)$ is a pure
function of the pass, the point and the lane's seed (with a different salt $s$ for value and chance), so the same pass
always varies the same way, and \ui{neu würfeln} changes only the seed.

\paragraph{Gestures through the room (\fn{autoPipeTime}).} A lane on a per-pipe target is not heard by all pipes at
once: a pipe at $\mathbf x$ (from its direction and distance, $x = -d\sin\varphi\cos\vartheta$,
$y = d\cos\varphi\cos\vartheta$, $z = d\sin\vartheta$) reads the lane at the earlier time
\begin{equation}
  t_\mathbf{x} = t - \frac{\lVert\mathbf x - \mathbf o\rVert}{v_w}\qquad\text{or, inwards,}\qquad
  t_\mathbf{x} = t - \frac{\max(0,\,25 - \lVert\mathbf x-\mathbf o\rVert)}{v_w},
\end{equation}
with the origin $\mathbf o$ and the speed of the front $v_w$ in metres per beat: the gesture spreads like a wave from
the origin (or runs in from a sphere of \SI{25}{m}). Its value $v$ acts on the target as the offset $2v-1$, a routing at
full amount.

\paragraph{Thinning recorded curves (\fn{autoThin}).} A recorded movement produces a point every few milliseconds. It is
reduced with the Ramer--Douglas--Peucker algorithm \citep{douglas1973}: keep the two end points; find the point farthest
from the straight line between them, $d = |p_x\Delta y - p_y\Delta x|/\sqrt{\Delta x^2+\Delta y^2}$ (time scaled into
value units); if it is farther than the tolerance, keep it and repeat on both halves. The result is the smallest set of
points (in this greedy sense) that stays within the tolerance of the original everywhere.
